
\subsection{滤过 $A_\infty$ 同伦的构造 I：捷径}
\label{subsec:7-2-11}

\begin{proof}[命题~\ref{prop:7-2-93} 的证明]
我们将直接由命题~\ref{prop:7-2-126}
推出命题~\ref{prop:7-2-93}。
取 $L=L'$，并令 $\psi$ 为恒等映射。
于是有包含 $\mathcal X(L,\lambda')$ 的
可数集 $\mathcal X(L,(\lambda,\lambda'))$。
这个集合就是第~\ref{subsec:7-2-9} 小节中
记为 $\mathcal X_3(L',(\lambda,\lambda'))$ 的集合；
在开始证明命题~\ref{prop:7-2-132} 之前，
这里暂改记号以免混淆。
$\mathcal X(L,(\lambda,\lambda'))$
生成 $C(L,(\lambda,\lambda'))$。
此外还有滤过
$A_{n'(\lambda),K'(\lambda)}$ 同态
\[
\mathfrak h^{(\lambda,\lambda'),-}\colon
C(L;\lambda)\longrightarrow C(L,(\lambda,\lambda')) .
\]
将 $\mathfrak h^{(\lambda,\lambda'),-}$
与包含映射
$C(L;\lambda')\subset C(L,(\lambda,\lambda'))$
的同伦逆复合，便得到所需的滤过
$A_{n'(\lambda),K'(\lambda)}$ 同伦等价
\[
\mathfrak h^{(\lambda,\lambda')}\colon
C(L,\lambda)\longrightarrow C(L,\lambda').
\]
\end{proof}

\begin{proof}[命题~\ref{prop:7-2-132} 的证明]
由上面证明的命题~\ref{prop:7-2-93}，
我们得到滤过
$A_{n'(\lambda),K'(\lambda)}$ 同伦等价
\[
\begin{aligned}
\mathfrak h^{1,(\lambda,\lambda')}&\colon
C_1(L,\lambda)\longrightarrow C_1(L,\lambda'),\\
\mathfrak h^{2,(\lambda,\lambda')}&\colon
C_2(L',\lambda)\longrightarrow C_2(L',\lambda').
\end{aligned}
\]
为证明
$\mathfrak g^{\lambda'}\circ
\mathfrak h^{1,(\lambda,\lambda')}$
与
$\mathfrak h^{2,(\lambda,\lambda')}\circ
\mathfrak g^\lambda$
在 $A_{n''(\lambda),K''(\lambda)}$ 意义下同伦，
如下应用命题~\ref{prop:7-2-160}
与命题~\ref{prop:7-2-162}。
注意
$\mathfrak h^{1,(\lambda,\lambda')}$、
$\mathfrak h^{2,(\lambda,\lambda')}$
是将命题~\ref{prop:7-2-126}
应用于恒等映射得到的滤过
$A_{n'(\lambda),K'(\lambda)}$ 同态；
$\mathfrak f^{(\lambda,\lambda)}$
则由该命题应用于 $\psi$ 得到，
而 $\mathfrak g^\lambda$ 是它与
$C_2(L',\lambda)\to C_3(L',(\lambda,\lambda))$
的同伦逆的复合。
显然 $\mathrm{id}\circ\psi=\psi\circ\mathrm{id}$。
因此，由定理 4.6.44 与
定理~\ref{thm:4-6-25} 的 $A_{n,K}$ 版本，
可知上述两个复合在
$A_{n''(\lambda),K''(\lambda)}$ 意义下同伦。
定理~\ref{thm:4-6-25} 的 $A_{n,K}$ 版本
已在第~\ref{subsec:7-2-10} 小节证明；
第 4.6.3 小节给出的定理 4.6.44 的证明，
在目前阶段也蕴含其 $A_{n,K}$ 版本。
\end{proof}

最后详细说明怎样从定理~\ref{thm:4-6-25}
的 $A_{n,K}$ 版本得到 $A_\infty$ 版本。

先给出一个适用于 $R$ 为有限域或 $\mathbb C$ 的捷径。
本书的大多数应用只需这两种情形。
下一引理中，设 $C_i$ 是
$\Lambda_{0,\mathrm{nov}}(R)$ 上的滤过
$A_\infty$ 代数。

\begin{lemma}{7.2.177}\label{lem:7-2-177}
设 $R$ 为有限域或 $\mathbb C$，
$C_1,C_2$ 为有隙的滤过 $A_\infty$ 代数。
假定各 $C_i$ 的 $R$-归约 $\overline C_i$
是有限生成自由 $R$ 模。
设 $\mathfrak g,\mathfrak g'$ 是从 $C_1$ 到 $C_2$
的有隙滤过 $A_\infty$ 同态。
若对所有 $(n,K)$，
$\mathfrak g$ 与 $\mathfrak g'$
均滤过 $A_{n,K}$ 同伦，
则二者有隙滤过 $A_\infty$ 同伦。
\end{lemma}

\par\noindent\textit{引理~\ref{lem:7-2-177} 的证明。}
采用定义--引理 4.2.30 中的模型
$\mathfrak C_2=C_2^{[0,1]}$。
由于 $\overline C_2$ 是有限维 $R$ 向量空间，
$\overline{\mathfrak C}_2$ 也为有限维
$R$ 向量空间；另一个模型
$\operatorname{Poly}([0,1],C_2)$
则不具有这一性质。
令 $\mathfrak{HOM}_{n,K}(C_1,\mathfrak C_2)$
为所有滤过 $A_{n,K}$ 同态
$C_1\to\mathfrak C_2$ 的集合。
它包含于有限多个如下 $R$ 向量空间的直和：
\[
\operatorname{Hom}_R
\bigl(B_k\overline C_1[1],
\overline{\mathfrak C}_2[1]\bigr).
\]
故它包含于一个有限维 $R$ 仿射空间。
滤过 $A_{n,K}$ 同态的定义方程是多项式型的，
因此
$\mathfrak{HOM}_{n,K}(C_1,\mathfrak C_2)$
可视为 $R$ 上的仿射代数簇。

对 $(n,K)<(n',K')$，显然的限制给出
代数簇态射
\[
\operatorname{Res}_{(n',K'),(n,K)}\colon
\mathfrak{HOM}_{n',K'}(C_1,\mathfrak C_2)
\longrightarrow
\mathfrak{HOM}_{n,K}(C_1,\mathfrak C_2).
\]
类似地定义 $\mathfrak{HOM}_{n,K}(C_1,C_2)$。
于是 $\operatorname{Eval}_{s=s_0}$ 给出态射
\[
(\operatorname{Eval}_{s=s_0})_{(n,K),*}\colon
\mathfrak{HOM}_{n,K}(C_1,\mathfrak C_2)
\longrightarrow
\mathfrak{HOM}_{n,K}(C_1,C_2).
\]
令
\begin{equation}
\begin{aligned}
V(n,K)=\bigl\{\mathfrak h\in
\mathfrak{HOM}_{n,K}(C_1,\mathfrak C_2)
\mid{}&
(\operatorname{Eval}_{s=0})_{(n,K),*}(\mathfrak h)
=\mathfrak g,\\
&
(\operatorname{Eval}_{s=1})_{(n,K),*}(\mathfrak h)
=\mathfrak g'\bigr\}.
\end{aligned}
\tag{7.2.178}\label{eq:7-2-178}
\end{equation}
由假设，每个 $V(n,K)$ 均非空。
它也是 $R$ 上的仿射代数簇，
且限制给出态射
\[
\operatorname{Res}_{(n',K'),(n,K)}\colon
V(n',K')\longrightarrow V(n,K).
\]
只须证明逆极限
\begin{equation}
\varprojlim V(n',K')
\tag{7.2.179}\label{eq:7-2-179}
\end{equation}
非空。若 $R$ 为有限域，则 $V(n,K)$
是有限集，而指标对 $(n,K)$ 构成有向集，
结论立即成立。

现在考虑 $R=\mathbb C$ 的情形。

\begin{statement}{子引理}{7.2.180}\label{sublem:7-2-180}\itshape
设 $R=\mathbb C$，并给定正整数 $I$。
则对每个 $i\le I$，存在
$\mathfrak h_i\in V(n_i,K_i)$，使所有
$i<i'\le I$ 都满足
\[
\operatorname{Res}_{(n(i'),K(i')),(n(i),K(i))}
(\mathfrak h_{i'})=\mathfrak h_i,
\]
并且
\[
\mathfrak h_i\in
\bigcap_{j>i}\operatorname{Im}
\bigl(\operatorname{Res}_{(n(j),K(j)),(n(i),K(i))}\bigr).
\]
\end{statement}

\begin{proof}
先证 $I=1$。令
\[
Z_j=\operatorname{Im}
\bigl(\operatorname{Res}_{(n(j),K(j)),(n(1),K(1))}\bigr).
\]
由假设，各 $Z_j$ 非空，且显然随 $j$ 递减。
我们要证明 $\bigcap_jZ_j$ 非空。
依 Chevalley 定理，每个 $Z_j$ 都是
\emph{可构造集}\liftIndex{I265}；
见 \textbf{[Mat70]} 第 42 页。
这里“可构造”意为它是关于扎里斯基拓扑的
有限个局部闭子集之并。
其维数不增，所以当 $j\ge j_0$、
$j_0$ 充分大时稳定为某个常数 $d$。
维数为 $d$ 的不可约分支数也不增，
故可假定它已稳定。
在 $Z_{j_0}$ 中取一个 $d$ 维不可约分支 $C$，
并对 $j>j_0$ 取 $Z_j$ 中相应的分支 $C_j$。
由于 $\dim C=\dim C_j$ 且 $C\supseteq C_j$，
$C\setminus C_j$ 在复解析拓扑下无处稠密。
因此，贝尔纲定理表明
$\bigcup_{j>j_0}(C\setminus C_j)$
为第一纲集，特别不能覆盖 $C$。
\translatornote{原书称上述可数并仍“无处稠密”；
一般只能断言它是第一纲集，非空交的结论
仍由贝尔纲定理得到。}
于是 $\bigcap_{j\ge j_0}C_j$ 非空。
按定义它包含于 $\bigcap_jZ_j$，
故可取
$\mathfrak h_1\in\bigcap_jZ_j$。

现假定子引理已对 $I$ 成立。
对 $j>I$，令
\[
W_{I,j}=\bigl\{\mathfrak h\in V(n(j),K(j))
\mid
\operatorname{Res}_{(n(j),K(j)),(n(I),K(I))}
(\mathfrak h)=\mathfrak h_I\bigr\}.
\]
同理可知每个 $W_{I,j}$ 非空。
又因为
$\operatorname{Res}_{(n(j),K(j)),(n(I+1),K(I+1))}
(W_{I,j})$
是一列递减的可构造集，
可用同样方法证明其交非空，
从而取到 $\mathfrak h_{I+1}$。
\end{proof}

这就完成了引理~\ref{lem:7-2-177} 的证明。\hfill$\square$

现在可在 $R=\mathbb C$ 或有限域的情形下
完成定理~\ref{thm:4-6-25} 的证明。
回顾已经得到的结果：
$C_1(L,\lambda_0)$ 与 $C_2(L',\lambda_0)$
上的滤过 $A_\infty$ 结构
已在第~\ref{subsec:7-2-8} 小节构造；
其中暂缓的命题~\ref{prop:7-2-93}
已在本小节证明。
它们之间的滤过 $A_\infty$ 同态
\[
\mathfrak g_{s=0}^{\psi_0,\{J_\rho^{s=0}\}_\rho,
\mathfrak s_{s=0}^{\operatorname{top}(\rho)}},
\qquad
\mathfrak g_{s=1}^{\psi_1,\{J_\rho^{s=1}\}_\rho,
\mathfrak s_{s=1}^{\operatorname{top}(\rho)}}
\]
已在第~\ref{subsec:7-2-9} 小节构造；
其中暂缓的命题~\ref{prop:7-2-132}
也已在本小节证明。
第~\ref{subsec:7-2-10} 小节证明了：
对任意 $n,K$，这两个同态
在 $A_{n,K}$ 意义下同伦。

由定理 5.4.1，
$C_1(L,\lambda_0)$ 与 $C_2(L',\lambda_0)$
分别滤过 $A_\infty$ 同伦等价于
$H(L,\Lambda_{0,\mathrm{nov}})$ 与
$H(L',\Lambda_{0,\mathrm{nov}})$
上的滤过 $A_\infty$ 代数。
后两者均为有限生成。
故可应用引理~\ref{lem:7-2-177}，
得到上述两个 $\mathfrak g$ 同态
滤过 $A_\infty$ 同伦。
定理~\ref{thm:4-6-25}
在 $R=\mathbb C$ 或有限域的情形下得证。

\begin{remark}{7.2.181}\label{rem:7-2-181}
遗憾的是，上述引理~\ref{lem:7-2-177}
的证明不适用于一般系数环 $R$。
例如，$\mathbb Q$ 上非空代数簇的有理点
所成的逆极限一般可能为空。
\translatornote{原书写“几何点”，
但随后的集合 $X_n\subset\mathbb Q^n$
讨论的是 $\mathbb Q$-有理点；
这里按例子的实际论证范围表述。}
考虑集合列
\[
X_n=\bigl\{(x_1,\ldots,x_n)\in\mathbb Q^n
\mid x_i=1+x_{i+1}^2,
\ i=1,\ldots,n-1\bigr\}.
\]
遗忘 $x_{n+1}$ 给出投影
$X_{n+1}\to X_n$。
当 $n\ge2$ 时，
若 $(x_1,\ldots,x_n)\in X_n$，
归纳可得 $x_1\ge n-1$。
\translatornote{原书未限定 $n\ge2$；
$X_1=\mathbb Q$，故其所称
“$x_1\ge n-1$”在 $n=1$ 时不成立。}
于是 $\varprojlim X_n=\varnothing$。
作者不知道引理~\ref{lem:7-2-177}
对一般 $R$ 是否成立。

这一困难与逆极限函子 $\varprojlim$
在各种范畴中并非正合函子有关。
这一点在代数拓扑中已见于
\textbf{[Mil62]}。
\end{remark}
